% GENERATED FILE. Edit canonical lessons, then run npm run generate:books.
% canonical-source-hash: fnv1a64:4333ad5087ade2cf
% canonical-lessons: FR-C28-cafe, FR-C28-the, FR-C28-lait, FR-C28-sucre

\chapter{Coffee, Tea, Milk, Sugar}
\label{ch:fr-coffee-tea-milk-sugar}
\hlchaptermodality{28}

\begin{chapteropening}
\textbf{By the end of this chapter:} \emph{I can ask for coffee, tea, milk or sugar in French, and trace each word to the road it travelled to get there — one unbroken, one forked, one never borrowed, one from Sanskrit.}

Ask for sugar, trace sucre to Sanskrit śárkarā through Arabic sukkar, and set the chapter's four masculine drink-words beside l'eau's lone feminine exception. Reaches back past the chapter to Chapter 18's si, the answer that contradicts a negative, and to the langue d'oïl / langue d'oc split behind oui itself.
\end{chapteropening}

\section[le café]{le café — coffee, one word the whole road agrees on}

\label{lesson:FR-C28-cafe}



\begin{quote}
You already know how to ask nicely — \textbf{s'il vous plaît}. Now something worth asking for: the first of four everyday words for what's on the table.
\end{quote}

\subsection*{You'll want to know: the word}
\textbf{le café} — coffee. \textbf{Masculine}.

\begin{quote}
\textbf{Un café, s'il vous plaît.} — \textquotedblleft{}A coffee, please.\textquotedblright{}
\end{quote}

\emph{Café} also names the \textbf{place} you drink it, exactly the way English \emph{cafe} does — one word, two jobs, in both languages at once.

\begin{cousinweb}
\textbf{café} comes from Ottoman Turkish \textbf{kahve}, which comes from Arabic \textbf{qahwa}. Arabic lexicographers record \emph{qahwa} as an older word for \textbf{wine} — it only attached to the coffee bean once the drink spread out of Yemen, in the fifteenth century.

From Ottoman \emph{kahve}, the word crossed into Europe through Italian \textbf{caffè} — and stopped there. French \textbf{café}, English \textbf{coffee}, German \textbf{Kaffee}, Italian \textbf{caffè} itself: every one of these is the \emph{same} borrowing, taken at a different stop on a \emph{single} road out of the Ottoman Empire. Nobody translated it; everybody just repeated it.
\end{cousinweb}

\subsection*{Your turn}
Say these aloud:

\begin{itemize}
\raggedright
  \item \textquotedblleft{}le café\textquotedblright{} — masculine
  \item \textquotedblleft{}un café, s'il vous plaît\textquotedblright{}
  \item the road — \textquotedblleft{}qahwa, kahve, caffè, café\textquotedblright{} — one word, four stops
  \item the surprise — \textquotedblleft{}qahwa once meant \textbf{wine}\textquotedblright{}
\end{itemize}

\begin{tcolorbox}[breakable,colback=teal!4,colframe=teal!35!black,title={Before you move on}]
Say \textquotedblleft{}coffee, please.\textquotedblright{} (\emph{\textbf{Un café, s'il vous plaît.}} — masculine.) What Arabic word is \emph{café} from, and what did it mean before coffee existed? (\emph{\textbf{Qahwa}} — \textquotedblleft{}\textbf{wine}.\textquotedblright{}) Name the road it travelled. (\textbf{Arabic qahwa $\to$ Ottoman kahve $\to$ Italian caffè $\to$ French café}, with English \emph{coffee} and German \emph{Kaffee} the same borrowing at other stops.) Is \emph{s'il vous plaît} the right register here? (\textbf{Yes} — spoken, to someone you don't know well.)
\end{tcolorbox}
\section[le thé]{le thé — tea, the word that took the other road}

\label{lesson:FR-C28-the}



\begin{quote}
You just saw \emph{café} travel Europe as one word, unchanged at every stop. Tea took the opposite kind of trip — it \textbf{split}.
\end{quote}

\subsection*{You'll want to know: the word}
\textbf{le thé} — tea. \textbf{Masculine}.

\begin{quote}
\textbf{Un thé, s'il vous plaît.} — \textquotedblleft{}A tea, please.\textquotedblright{}
\end{quote}

\begin{cousinweb}
China writes \textquotedblleft{}tea\textquotedblright{} with one character everywhere, but doesn't say it the same way everywhere. In Mandarin it's \textbf{chá}. In the \textbf{Hokkien} of the south-eastern coast — the Amoy dialect — it's \textbf{tê}.

Traders who bought the leaf overland or through Mandarin-speaking ports carried home \textbf{chá}: it travelled into Persian, into Arabic, into Hindi. Dutch traders bought their tea specifically at \textbf{Amoy}, in Hokkien territory, and carried home \textbf{tê} instead — which became Dutch \textbf{thee}, then French \textbf{thé}, English \textbf{tea}, German \textbf{Tee}.

So \emph{café} and \emph{thé} make a matched pair with opposite stories: \emph{café} is one word repeated unchanged at every stop; \emph{thé} is the \emph{same original word}, worn into two different shapes depending which Chinese port a trader happened to deal with.
\end{cousinweb}

\subsection*{Your turn}
Say these aloud:

\begin{itemize}
\raggedright
  \item \textquotedblleft{}le thé\textquotedblright{} — masculine
  \item \textquotedblleft{}un café, un thé\textquotedblright{} — both masculine, both borrowed
  \item the fork — \textquotedblleft{}chá overland, tê by Dutch ship\textquotedblright{}
  \item the pair — \textquotedblleft{}café: one road. thé: two.\textquotedblright{}
\end{itemize}

\begin{tcolorbox}[breakable,colback=teal!4,colframe=teal!35!black,title={Before you move on}]
Say \textquotedblleft{}tea, please.\textquotedblright{} (\emph{\textbf{Un thé, s'il vous plaît.}} — masculine.) What Chinese dialect gave French \emph{thé} its shape, and by what route? (\textbf{Hokkien tê}, carried by \textbf{Dutch} ships from the port of \textbf{Amoy}.) What's the overland word that \emph{chá}-family languages use instead? (\textbf{Chá}, into \textbf{Persian, Arabic, Hindi}.) How is \emph{thé}'s story different from \emph{café}'s? (\emph{Café} stayed \textbf{one word} everywhere; \emph{thé} \textbf{split in two} depending on the port.)
\end{tcolorbox}
\section[le lait]{le lait — milk, the word French never had to borrow}

\label{lesson:FR-C28-lait}



\begin{quote}
\emph{Café} and \emph{thé} both crossed an ocean to reach French. This word never left home.
\end{quote}

\subsection*{You'll want to know: the word}
\textbf{le lait} — milk. \textbf{Masculine}, like \emph{café} and \emph{thé}.

\begin{quote}
\textbf{Du lait, s'il vous plaît.} — \textquotedblleft{}Some milk, please.\textquotedblright{}
\end{quote}

\begin{cousinweb}
\textbf{lait} comes from Latin \emph{\textbf{lac}} (stem \emph{\textbf{lact-}}) — a word Latin-speakers in Gaul simply kept saying, generation after generation, until it wore down into \emph{lait}. No port, no trade route: this one was \textbf{inherited}, not borrowed, unlike \emph{café}'s single road or \emph{thé}'s forked one.

English shows the same split \emph{tête} showed for \textquotedblleft{}head\textquotedblright{}: the \textbf{inherited}, everyday English word for this drink is \textbf{milk} — and it has \textbf{nothing to do with} \emph{lac}. \emph{Milk} descends from a completely different Proto-Indo-European root, one meaning \textquotedblleft{}to milk a cow.\textquotedblright{} English only meets \emph{lac}'s stem \emph{lact-} in \textbf{learned} words borrowed straight from written Latin much later: \textbf{lactic}, \textbf{lactose}, \textbf{lactate}. So French \emph{lait} and English \emph{lactic} are close cousins; French \emph{lait} and English \emph{milk} are not related at all.
\end{cousinweb}

\subsection*{Your turn}
Say these aloud:

\begin{itemize}
\raggedright
  \item \textquotedblleft{}le lait\textquotedblright{} — masculine
  \item \textquotedblleft{}un café, un thé, du lait\textquotedblright{} — three drinks, one gender
  \item the cousins — \textquotedblleft{}lait, lactic, lactose\textquotedblright{}
  \item the non-cousins — \textquotedblleft{}lait and milk share no root at all\textquotedblright{}
\end{itemize}

\begin{tcolorbox}[breakable,colback=teal!4,colframe=teal!35!black,title={Before you move on}]
Say \textquotedblleft{}milk, please.\textquotedblright{} (\emph{\textbf{Du lait, s'il vous plaît.}} — masculine.) Is \emph{lait} borrowed, like \emph{café} and \emph{thé}, or inherited straight from Latin? (\textbf{Inherited} — from \emph{lac}, stem \emph{lact-}.) Which English words are close cousins of \emph{lait}? (\textbf{Lactic, lactose, lactate} — learned borrowings of the same stem.) Is English \emph{milk} one of those cousins? (\textbf{No} — a completely different, unrelated root.)
\end{tcolorbox}
\section[le sucre]{le sucre — sugar, the longest road of the four}

\label{lesson:FR-C28-sucre}



\begin{quote}
Coffee took one road, tea forked into two, milk never left home. Sugar's road is the longest of all — and it ends the chapter with a whole table's worth of masculine drink-words.
\end{quote}

\subsection*{You'll want to know: the word}
\textbf{le sucre} — sugar. \textbf{Masculine.}

\begin{quote}
\textbf{Du sucre, s'il vous plaît.} — \textquotedblleft{}Some sugar, please.\textquotedblright{}
\end{quote}

\begin{cousinweb}
\textbf{sucre} $\leftarrow$ Old French \textbf{çucre} $\leftarrow$ Arabic \textbf{sukkar} $\leftarrow$ Persian \textbf{shakar} $\leftarrow$ Sanskrit \textbf{śárkarā}, which first meant \textquotedblleft{}\textbf{grit, gravel}\textquotedblright{} — crystallized sugar looked like gravel to whoever named it first.

\emph{Sucre} passed through Arabic just like \emph{café} did, but through a \textbf{different} Arabic word (\emph{sukkar}, not \emph{qahwa}): a shared \textbf{language of transmission}, not a shared root. Where \emph{lait} is inherited straight from Latin, \emph{sucre} is borrowed — from \textbf{further back} than \emph{café} or \emph{thé}, all the way to Sanskrit.
\end{cousinweb}

\begin{grammarlens}[title={Grammar Lens: a lean, and its one loud exception}]
Four drink-words, four \textbf{masculine} articles: \emph{le café, le thé, le lait, le sucre}. Not a rule to trust — a lean: French tends to make plain substances masculine. This lean has a loud exception: \textbf{l'eau}, \textquotedblleft{}water,\textquotedblright{} is \textbf{feminine}. You still learn each word with its article, the promise this book has kept since \emph{la main}.
\end{grammarlens}

\begin{culture}
Chapter 18's \textbf{si} — the \textquotedblleft{}yes\textquotedblright{} for contradicting a negative — fits this table exactly:

\begin{quote}
\textbf{Vous ne voulez pas de sucre ?} — \emph{\textbf{Si, un peu.}} (\textquotedblleft{}You don't want sugar? — Yes, a little.\textquotedblright{})
\end{quote}

Plain \textbf{oui} here would agree you \emph{don't} want it; \emph{si} is what pushes back. That same \emph{oui} once mattered enough to \textbf{split a country} — \emph{langue d'oïl} in the north, \emph{langue d'oc} in the south.
\end{culture}

\subsection*{Your turn}
Say these aloud:

\begin{itemize}
\raggedright
  \item \textquotedblleft{}le sucre\textquotedblright{} — masculine
  \item all four — \textquotedblleft{}le café, le thé, le lait, le sucre\textquotedblright{}
  \item the exception — \textquotedblleft{}l'eau\textquotedblright{} — feminine, against the lean
  \item the contradiction — \textquotedblleft{}Vous ne voulez pas de sucre ? — Si !\textquotedblright{}
\end{itemize}

\begin{tcolorbox}[breakable,colback=teal!4,colframe=teal!35!black,title={Before you move on}]
Say \textquotedblleft{}sugar, please.\textquotedblright{} (\emph{\textbf{Du sucre, s'il vous plaît.}} — masculine.) Trace \emph{sucre} back to Sanskrit. (\textbf{Sukkar $\leftarrow$ shakar $\leftarrow$ śárkarā}, \textquotedblleft{}\textbf{grit, gravel}.\textquotedblright{}) Do \emph{café} and \emph{sucre} share a root? (\textbf{No} — a language of transmission, Arabic, not a root; \emph{qahwa} and \emph{sukkar} are different words.) What is the lean about French drink-words, and what breaks it? (\textbf{Masculine by default} — broken by \textbf{l'eau}, feminine.) When do you answer \textbf{si} instead of \textbf{oui}? (\textbf{Contradicting a negative} — the same word-pair that once \textbf{split France} into \emph{langue d'oïl} and \emph{langue d'oc}.)
\end{tcolorbox}
